\section{Old-Task Loss Preservation}

Since $\ell_A$ is our best practical surrogate on performance on distribution $A$, a natural choice is to use $\ell_A$ itself to define preservation. We say that the model $f_\theta$ preserved the solution if it preserved the old loss value for distribution $A$. This leads to:
\[
q(a,\theta)
=
\ell_A(f_\theta(x))-\ell_A(f_{\theta_A^\star}(x)).
\]
Since preservation means $q(a,\theta)=0$, take $\mathcal C=\{0\}$. The corresponding unscaled proximal loss is
\[
p(a,\theta)
=
\operatorname{dist}(q(a,\theta),\mathcal C)^2
=
q(a,\theta)^2.
\]
To obtain a low-order Taylor approximation, use the perturbation around the old solution $\delta=\theta-\theta_A^\star$ in the first-order expansion of $q$:
\[
q(a,\theta_A^\star+\delta)
\approx
g^\top\delta,
\]
where
\[
g
=
\nabla_\theta \ell_A(f_{\theta_A^\star}(a)).
\]

Equivalently, writing
\[
z_A=f_{\theta_A^\star}(a)
\qquad
r_i
=
\nabla_z \ell_A(z)\big|_{z=f_{\theta_A^\star}(a)},
\qquad
J_i
=
\nabla_\theta f_{\theta_A^\star}(a),
\]
the chain rule gives
\[
g_i
=
J_i^\top r_i.
\]

Substituting this expansion into the preservation penalty gives
\[
\frac{\lambda}{2N_A}\sum_{i=1}^{N_A} q(a_i,\theta)^2
\;\approx\;
\frac{\lambda}{2N_A}\sum_{i=1}^{N_A} (g_i^\top\delta)^2.
\]
Collecting the gradient outer products yields
\[
\frac{\lambda}{2N_A}\sum_{i=1}^{N_A} (g_i^\top\delta)^2
=
\frac{\lambda}{2}\,
\delta^\top
\left(
\frac{1}{N_A}\sum_{i=1}^{N_A} g_i g_i^\top
\right)
\delta
=
\frac{\lambda}{2}\,
\delta^\top
M_1
\delta.
\]

Therefore the resulting matrix is
\[
M_1
=
\frac{1}{N_A}\sum_{i=1}^{N_A} g_i g_i^\top
=
\frac{1}{N_A}\sum_{i=1}^{N_A} J_i^\top r_i r_i^\top J_i.
\]
We call this an importance matrix as it encodes how important it is to preserve a direction in parameter space.
